% Status: history (G2). Superseded by notes/channels_system.tex, notes/dsl_updates.md + castlegen/channels. Kept for the record.
\documentclass[11pt]{article}
\usepackage[margin=1in]{geometry}
\usepackage{amsmath,amssymb}
\usepackage{booktabs}
\usepackage{dsfont}
\usepackage{algorithm}
\usepackage{algpseudocode}

\newcommand{\Z}{\mathbb{Z}}
\newcommand{\R}{\mathbb{R}}
\newcommand{\1}{\mathds{1}}
\newcommand{\Dom}{\mathcal{D}}
\newcommand{\Chan}{\mathcal{C}}
\newcommand{\Cand}{\mathcal{Q}}
\newcommand{\refine}{\mathcal{R}}

\title{Channels: one interface for latents and tiles}
\author{Hierwave notes (working document)}
\date{}

\begin{document}
\sloppy
\maketitle

\noindent Purpose: fix a notation in which every variable the sampler touches
(exemplar coordinates, promises, object codes, macro rooms, certificate fields,
tiles) is an instance of one thing, a \emph{channel}, so that adding a channel
adds tables and declared couplings and nothing else. Companion to
\texttt{coarse\_to\_fine.tex} (target, generator, promise machinery), which this
document does not repeat; \texttt{dsl.md} \S4--5 and \texttt{macro\_objects.md}
are the prose sources.

Sections~\ref{sec:notation}--\ref{sec:existing} state what exists in the
notation. Section~\ref{sec:open} lists what is not yet settled; it is the agenda.

%-----------------------------------------------------------------------------
\section{Notation}
\label{sec:notation}

\paragraph{Lattices.} Levels are cell sides $h \in \{h_{\text{top}}, \dots, 2, 1\}$
in tile units. $\Lambda_h = (h\Z)^2$ is the level-$h$ lattice; a \emph{cell}
$p \in \Lambda_h$ covers the $h \times h$ tiles $[p, p + h)^2$. $\mathrm{par}(p)
\in \Lambda_{2h}$ is the parent of $p$, $\mathrm{ch}(p) \subset \Lambda_{h/2}$ its
four children, $N_r(p) \subset \Lambda_h$ the same-level cells within Chebyshev
distance $r$ (in cells). A chunked world is a finite window of each $\Lambda_h$
with a halo; nothing below depends on the window.

\paragraph{Channels.} A channel $c \in \Chan$ lives at one level $h_c$ and has a
finite domain $\Dom_c$. Its state is a field $z_c : \Lambda_{h_c} \to \Dom_c$.
The channels at level $h$ are $\Chan_h$, and $z_h = (z_c)_{c \in \Chan_h}$ is the
level's state; $z_h(p) \in \prod_{c \in \Chan_h} \Dom_c$ is one cell's state. Tiles
are the channel $t$ with $h_t = 1$ and $\Dom_t = \mathcal T$. A channel is
\emph{small} if $|\Dom_c|$ is enumerable per update (promises, codes, certificate
fields; tens to a few thousand) and \emph{large} otherwise (exemplar coordinates,
$m^2$; macro rooms, $K \cdot \mathrm{BS}^2$).

\paragraph{Abstractions.} A \emph{derived} channel has an abstraction
\[
  \alpha_c : \mathcal T^{\,h_c \times h_c} \to \Dom_c,
  \qquad z_c(p) = \alpha_c\big(x|_{p}\big)\ \text{in the target},
\]
a deterministic summary of the tiles under the cell. It is \emph{compositional}
if there is $m_c$ with $\alpha_c(x|_p) = m_c\big((\alpha_c(x|_q))_{q \in \mathrm{ch}(p)},
\text{seams}\big)$, which gives an exact parent--child relation for free. A
\emph{free} channel has no $\alpha$ (exemplar coordinates: a statement about the
tiles, compared softly, not a function of them).

\paragraph{Factors.} Every term of the energy is one of three kinds, each either
soft (a table $E \in \R \cup \{+\infty\}$) or hard (an indicator; written as
$E \in \{0, +\infty\}$):
\begin{center}
\small
\begin{tabular}{@{}llll@{}}
\toprule
kind & arguments & scope & examples \\
\midrule
neighbour $E^{\mathrm{nb}}_{c,c'}$ & $z_c(p),\, z_{c'}(q)$, $q \in N_{r}(p)$ & $2r+1$ cells & $E_c$, pair tables, compat, overlap \\
cell $E^{\mathrm{cell}}_{c,c'}$ & $z_c(p),\, z_{c'}(p)$ & one cell & $\mathrm{cost}(u,v)$, $-\log z$, $t = \mathrm{tmpl}[o]$, $V_p$ \\
parent $E^{\mathrm{par}}_{c,c'}$ & $z_c(p),\, z_{c'}(\mathrm{par}\,p)$ & two levels & $E_{\mathrm{par}}$, refinement, $\nu$ \\
\bottomrule
\end{tabular}
\end{center}
Hard neighbour factors of radius $r > 1$ exist (a room's footprint reaches
$\mathrm{RAD} = 3$ blocks); so far every soft neighbour factor has $r \le 2$.

%-----------------------------------------------------------------------------
\section{Target and generator in channel terms}

As in \texttt{coarse\_to\_fine.tex}: the \emph{target} is undirected on tiles,
$p_G(x) \propto e^{-E_p(x)/T}\,\1[G(x)]$, and derived channels are its
push-forward through $\alpha$. The \emph{generator} is a chain graph whose chain
components are the levels, directed coarse to fine:
\begin{gather}
  g(z) = g_{\text{top}}(z_{h_{\text{top}}})\prod_{h < h_{\text{top}}} g_h(z_h \mid z_{2h}), \notag \\
  g_h(z_h \mid z_{2h}) = \frac{1}{Z_h(z_{2h})}
  \exp\!\Big(-\tfrac{1}{T}\big[E^{\mathrm{par}}_h(z_h, z_{2h}) + E^{\mathrm{nb}}_h(z_h) + E^{\mathrm{cell}}_h(z_h)\big]\Big),
  \label{eq:level}
\end{gather}
with the tile level $h = 1$ included as the last component. Within a level all
factors are undirected over the cells and channels of that level; between levels
every factor is directed downward. Local normalisation is what buys locality,
fittable tables and ancestral sampling, at the price that nothing below ever
pushes back on a level above (the \texttt{coarse\_to\_fine} $\mathrm{TV}$ bound).
Everything in this document is about the \emph{within-level} and
\emph{parent} factors of one level and how they are sampled.

%-----------------------------------------------------------------------------
\section{The channel interface}
\label{sec:interface}

What a channel must expose for the sampler to treat it generically. Items marked
$\dagger$ are present in the code for at least one channel; the rest are
proposals. Section~\ref{sec:views} reduces this list to views and factors.

\begin{center}
\small
\begin{tabular}{@{}lll@{}}
\toprule
member & meaning & code today \\
\midrule
$h_c$, $\Dom_c$ $\dagger$ & level and domain & every channel \\
$\alpha_c$, $m_c$ $\dagger$ & abstraction and merge (derived channels) & \texttt{leaf}, \texttt{merge} (generic.py) \\
$\refine_c(z_c(\mathrm{par}\,p) \mid \text{halo})$ $\dagger$ & child values a parent allows (hard) & \texttt{refine}, \texttt{compat} \\
$E^{\mathrm{nb}}, E^{\mathrm{cell}}, E^{\mathrm{par}}$ $\dagger$ & factor tables, by kind; hard ones as masks & scattered \\
$\Cand_c(p \mid z_{-p})$ $\dagger$ & candidates for a large domain, from neighbours & \texttt{\_mh\_S}, macroobj A \\
$\iota_d(z_c)$ & the interface a value presents across side $d$ & ports, seam projections \\
$\mathrm{cost}_c$ & work per update (for a compile-time report) & -- \\
\bottomrule
\end{tabular}
\end{center}

\paragraph{The cell update.} Fix a cell $p$ at level $h$ and a block of channels
$B \subseteq \Chan_h$ updated jointly. The conditional of $z_B(p)$ given everything
else is
\begin{align}
  \pi_p(z_B) \;\propto\; \exp\!\Big(-\tfrac1T \Big[
    &\sum_{c \in B} \Big( E^{\mathrm{par}}_c\big(z_c, z(\mathrm{par}\,p)\big)
    + \sum_{c'}\sum_{q \in N_r(p)} E^{\mathrm{nb}}_{c,c'}\big(z_c, z_{c'}(q)\big)
    + \sum_{c' \notin B} E^{\mathrm{cell}}_{c,c'}\big(z_c, z_{c'}(p)\big)\Big) \notag \\
    &+ \sum_{c, c' \in B} E^{\mathrm{cell}}_{c,c'}(z_c, z_{c'})
  \Big]\Big),
  \label{eq:cell}
\end{align}
where the neighbour sum includes the cells that $p$'s value is read by, not only
the ones $p$ reads (symmetry of $E^{\mathrm{nb}}$). Two update kinds, chosen by
the block's domain size:
\begin{itemize}
  \item \textbf{exact} when $\prod_{c \in B}|\Dom_c|$ is enumerable: every candidate
        scored, hard factors as masks, a Gumbel-max or integer draw (bit-sliced
        when the factors factor over port groups, \texttt{port\_gibbs.md});
  \item \textbf{proposal} otherwise: a candidate multiset $\Cand(p)$ built from the
        neighbours' values (coherent shifts, parent prolongation, $k$ links,
        $\epsilon$ uniform), scored exactly within it, with the
        Metropolis--Hastings correction wherever $\Cand$ depends on $z_B(p)$ itself.
\end{itemize}
A large channel paired with a small one (coordinate $u$ with promise $v$) is
updated by marginalising the small one inside each candidate's score,
$\bar\pi(u) = e^{-\Delta E(u)} \sum_v e^{-\mathrm{cost}(u,v)}\,\1[v \text{ compatible}]$,
then drawing $v \mid u$ (\texttt{joint\_sampler.md}).

\paragraph{Correctness conditions.} The following are what each update needs; they
are the generic replacements for per-channel arguments.
\begin{enumerate}
  \item \emph{Colouring.} Cells updated in parallel share no factor: with neighbour
        radius $r$ and a certificate that reads witnesses at distance 2, same-colour
        cells are $\ge 2r+1$ apart ($r = 2$: spacing 5; the plus-shaped
        neighbourhood needs the 5-colouring $(x + 2y) \bmod 5$).
  \item \emph{Never empty.} Under a hard neighbour factor the current value must
        remain a candidate: true for pair factors under a 2-colouring, and for
        certificates when the dependant checks are reproduced exactly.
  \item \emph{Valid proposals.} A candidate set that depends only on the neighbours
        and is sampled exactly, plus an $\epsilon$ uniform fallback, is a valid
        move; a set that depends on the current value needs a symmetric graph or
        the MH ratio with the reverse set.
  \item \emph{Partition rule.} Restricting an exact update to a subset must offer the
        same subset whichever member is current (a partition of the domain), else
        moves out of it cannot be reversed.
  \item \emph{Rates are free.} Channels in different blocks may be updated at any
        relative rate; each step leaves the joint invariant. This is what lets a
        matching variable be re-drawn once per ten sweeps while a hard channel is
        swept every time.
\end{enumerate}

\paragraph{Honourability.} With top-down sampling and no revisiting, for every
derived channel and every parent value its constraints allow, in every halo its
constraints allow,
\[
  \refine_c\big(z_c(\mathrm{par}\,p) \mid \text{halo}\big) \neq \varnothing,
\]
and the level-$K$ grid must have a tile completion. For one channel this is the
realisability property of a promise language. For several hard channels on one
cell it is a joint condition on $\prod_c \refine_c$, which no single language
checks; it is where joint promises (support with connectivity) can fail.

%-----------------------------------------------------------------------------
\section{The existing channels in this notation}
\label{sec:existing}

\begin{center}
\small
\begin{tabular}{@{}lllll@{}}
\toprule
channel & $h$ & $\Dom$ & $\alpha$ & update \\
\midrule
coordinate $u_h$ & every & $\Z_m^2$ & free & proposal \\
promise $v_h$ & $\ge K$ & small & $\alpha$, $m$ & exact, chain blocks \\
object code $o$ & 1 & codes & cell of template & exact, with $t$ \\
certificate $(g,d)$ & 1 & $(G{+}2)(D{+}1)$ & -- & closed form, with $t$ \\
macro room $(k,\phi)$ & block & $1 + K\,\mathrm{BS}^2$ & none & exact heat bath \\
tile $t$ & 1 & $\mathcal T$ & identity & exact, bit-sliced \\
\bottomrule
\end{tabular}

\medskip
\begin{tabular}{@{}lll@{}}
\toprule
channel & hard factors & soft factors \\
\midrule
coordinate $u_h$ & object seam check (nb) & $E_c$ (nb), $E_{\mathrm{par}}$ (par), cost (cell) \\
promise $v_h$ & compat (nb), refine (par) & $u$, pair, parent tables; cost (cell) \\
object code $o$ & nbr rule (nb), $t = \mathrm{tmpl}[o]$ (cell) & -- \\
certificate $(g,d)$ & validity $V_p$ (nb, $r = 2$) & $\mu g + \delta d$ (cell) \\
macro room $(k,\phi)$ & overlap, attachment or contact (nb, $r = 3$) & $\mu_{\mathrm{fam}}$, $-\lambda\,\#$contacts (nb), $\nu$ (par) \\
tile $t$ & rules, codes, ports (nb, cell) & $E_{\mathrm{loc}}$ (nb), parent term (par), $\tau$ \\
\bottomrule
\end{tabular}
\end{center}

\paragraph{Where the macro channel does not fit.} It is sampled alone on its own
lattice (blocks of $\mathrm{BS} = 8$) with a pasted coarse labelling as a fixed
parent, with no tile level under it (rooms are rendered from templates), with
connectivity as a heuristic (a spanning forest recomputed per sweep, which reads
the whole window) rather than a certificate or a promise, and with doors chosen
at render by a global union-find. Its energy
\[
  E(\text{rooms}) = \sum_{\text{rooms}} \mu_{\mathrm{fam}(k)} - \lambda\, \#\text{contacts}
  + \nu\, \#\{\text{blocks off the parent}\} + (\text{connectivity terms})
\]
is a neighbour factor of radius 3 in blocks plus a parent factor, so it fits
Eq.~\eqref{eq:cell}; what is missing is the parent relation $\refine$ and the
abstraction that a coarser level would carry, and the object codes that would
tie it to a tile level.

\paragraph{The matching variable.} \texttt{macro\_objects.md} \S1 proposes
$p(x, z) \propto e^{-E_{\mathrm{hard}}(x) - E_{\mathrm{soft}}(x, z)}$ with $z$ an
exemplar coordinate per block or window. In this notation $z$ is a coordinate
channel at the block level with $E^{\mathrm{cell}}(x, z)$ the $O(1)$ agreement
table and $E^{\mathrm{nb}}(z, z')$ the coherence of neighbouring matches; the
observation is rule 5 above (rates are free), and the current pipeline is the
case where $z$ is drawn once.

%-----------------------------------------------------------------------------
\section{Objects as a block-level coordinate channel}
\label{sec:objects}

Proposal: there is no separate ``macro room'' domain. The block level $h = B$
($B = 8$) carries an ordinary coordinate channel $u_B : \Lambda_B \to \Z_m^2$,
at tile stride, so a window can sit at any tile offset (this is where the phase
$\phi$ comes from), and objects are read off it by an abstraction. The macro
coarse level already is this (coordinates per window, pasted); the proposal is
that the fine macro stage is too.

\paragraph{Two abstractions, two hard rules.} Let $o_E$ be the exemplar's object
code channel and $W(u) = [u, u+B)^2$ the window a coordinate names.
\begin{itemize}
  \item \textbf{Seam form.} $\alpha^{\mathrm{seam}}_d(u) = o_E|_{\partial_d W(u)}$, the
        codes on the window's side-$d$ edge. Hard neighbour factor between
        adjacent cells: the \texttt{nbr} rule holds across the seam, pair by pair.
        This is \texttt{objects.coarse\_term}. Exact and whole-or-absent at every
        level.
  \item \textbf{Ownership form.} $\alpha^{\mathrm{own}}(u) \in A = \{\varnothing\} \cup
        \{(k,\phi)\}$: the object whose centre cell lies in $W(u)$, if any, with
        footprint $F(k,\phi) \subset \Z^2$ relative to the cell. Hard neighbour
        factor over $N_{\mathrm{RAD}}(p)$: $F_p \cap F_q = \varnothing$ (overlap), plus
        attachment or contact as a function of $(\alpha(p), \alpha(q))$. This is
        \texttt{macroobj} with $(k,\phi)$ replaced by $\alpha^{\mathrm{own}}(u)$.
\end{itemize}

\paragraph{Why the seam form freezes.} Under the seam form a cell's update must
keep every partial object on its four seams, so an object straddling a seam can
only be exchanged for another instance of the same template at the same offset.
An object of side $s$ lies inside one cell of side $h \ge s$ at a fraction
$\big((h - s + 1)/h\big)^2$ of positions: at $h = s$ almost none, at $h = 2s$
one in four. So at the level where the object ``is one cell'' most instances
still straddle, and the level above does not help, since there the cell holds
several objects and the window move drags all of them. The macro runs that
worked are ownership-form runs, and the tile-level finding that single-site
updates cannot move a rigid object is the seam form at $h = 1$.

\paragraph{The ownership form as $(a, u)$.} Hard terms depend on $u$ only
through $\alpha^{\mathrm{own}}(u)$; soft terms ($E_c$, $E_{\mathrm{par}}$) depend on
$u$ fully. Split the channel into $a_p \in A$ and $u_p \in \Z_m^2$ with
\[
  E^{\mathrm{cell}}(a, u) = \nu\,\1\big[a \ne \alpha^{\mathrm{own}}(u)\big],
\]
hard factors on $a$ alone, soft factors on $u$ alone. Then
\begin{itemize}
  \item $a_p \mid \text{rest}$ is an exact heat bath over $A$ ($|A| \approx 1300$):
        the current block update, with $\nu$ as the only soft term;
  \item $u_p \mid \text{rest}$ is a coordinate proposal update (coherent shifts,
        parent, $k$ links), scored by $E_c + E_{\mathrm{par}} + \nu\,\1[\cdot]$, run at
        whatever rate is wanted (rule 5, Section~\ref{sec:interface}).
\end{itemize}
This is \texttt{macro\_objects.md} \S1 stated as a channel: $\nu$ is not a
heuristic pull toward a pasted parent but the cell coupling between two channels
of one level, and the current pipeline is the case where $u$ is drawn once. At
$\nu \to \infty$ the pair collapses to one channel $u$ with $a = \alpha(u)$.

\paragraph{Masking.} With $a_p$ owning a footprint that runs into neighbouring
cells, the coordinate $u_q$ of a neighbour still states what its whole window
looks like, including the cells now under $F_p$. The statement must be masked
there: $E_c(u_q)$ and the tile term $\tau$ are evaluated over
$W(u_q) \setminus \bigcup_{p} F_p$, and the tile level reads $t = \mathrm{tmpl}[\cdot]$
on footprints and the coordinate elsewhere. This is the fixed-mask case of
\texttt{dsl.md} \S6, with the mask a deterministic function of the $a$ channel.

\paragraph{What it gives for free.}
\begin{itemize}
  \item \textbf{Interlocking.} Replace $F_p \cap F_q = \varnothing$ by agreement on the
        overlap, $\mathrm{tmpl}_p = \mathrm{tmpl}_q$ on $F_p \cap F_q$: shared walls and
        nested pieces are the same radius-$\mathrm{RAD}$ factor with a different
        table. The seam form would need successor sets for this.
  \item \textbf{The parent.} $u_B$ has the usual coordinate parent $u_{2B}$ with
        $E_{\mathrm{par}}$; no hard parent relation is needed for objects, since
        $\varnothing$ is always admissible. A parent that \emph{promises} an object
        (a grammar leaf, a required room) is a hard $E^{\mathrm{par}}$ on $a$, and
        honourability is then the question whether the footprint fits the halo.
  \item \textbf{Connectivity.} Contacts are a function of $(a_p, a_q)$, so the
        contact graph is a hard-table neighbour relation on $a$; the port promise
        and a certificate over objects-as-nodes attach to $a$, not to $u$.
\end{itemize}

\paragraph{What it costs.} Two channels at the block level instead of one, the
mask bookkeeping in $E_c$ and $\tau$, and a coordinate proposal set that knows
about object instances (a candidate for $u_p$ with $a_p$ fixed should prefer
windows whose $\alpha^{\mathrm{own}}$ equals $a_p$: the instances of $(k, \phi)$ in
the exemplar, which is a precomputed list).

%-----------------------------------------------------------------------------
\section{Derived rules, and planning in the coordinate update}
\label{sec:plan}

\paragraph{Non-overlap is derived, not authored.} Under a correct parent
$\to$ tile relation, two objects claiming one cell with different templates
have no tile completion: the exact coupling $F(a) = -\log Z_{\text{tiles}}(a)$ is
$+\infty$ on overlapping footprints, so the overlap rule would be learned (the
$F = \infty$ case of self-play, \texttt{dsl.md} \S7). It would also be found by
the honourability checker: it is exactly the set of $a$-configurations whose
$\prod \refine$ is empty. Either route needs machinery that does not exist, and
the rule is known in closed form, so it is hard-coded; in the source taxonomy of
\texttt{dsl.md} \S5 it is a \emph{derived} factor written by hand, not an
authored one. The same holds for interlocking (agreement on the overlap) and for
``a promised contact must have a door-capable wall behind it''. The general
principle: a hard neighbour factor at level $h$ is legitimate whenever it is
implied by the hard rules at level $h/2$ through $\alpha$; hand-writing it is a
shortcut for the compiler, and the checker, when it exists, verifies the hand
version.

\paragraph{Connectivity belongs in the plan.} In the macro runs the expensive
part was merging components by single-block moves on $a$: with no unmatched
door no single move merges two components, and the forest heuristic worked by
dissolving and regrowing. The coordinate update is the natural place to decide
where rooms connect, because a window $W(u)$ is a piece of the exemplar and is
connected inside exactly as the exemplar is; only the seams need planning. Give
the coordinate a second abstraction, its \emph{interfaces},
\[
  \iota_d(u) \in \mathcal I, \qquad d \in \{N, E, S, W\},
\]
the door-capable openings of $W(u)$'s objects on side $d$ (from $o_E$ and the
template's port sides; a $B$-bit pattern, or coarser: the half-sides that open).
Then:
\begin{itemize}
  \item \textbf{Seam factor on $u$.} $E^{\mathrm{nb}}_{u,u}(u_p, u_q) =
        \kappa\, \delta\big(\iota_d(u_p), \iota_{\bar d}(u_q)\big)$ for adjacent $p, q$: a
        proposal for $u_p$ is scored by whether its openings line up with the
        neighbours'. Read from the exemplar at both ends, so an $O(1)$ table
        per candidate, like $\nu$. Soft with $\kappa$, or hard on sides a promise
        requires joined.
  \item \textbf{The plan is honoured by $a$.} A three-way factor
        $E\big(a_p, a_q, \iota_d(u_p)\big)$: wherever the plan opens side $d$, the
        objects $a_p, a_q$ must be in contact across it (a hard mask on the
        $a$-update, built from the neighbour's $a_q$ and the own $u_p$, so still
        local). This replaces the forest heuristic: $a$ never has to \emph{discover}
        a merge, it has to realise an opening the plan already placed.
  \item \textbf{The global half stays a promise.} Interfaces give local joins
        only. Reachability needs the window-level promise as now (root side,
        keys): the contract ``every object reaches a contact across the root
        side'' is a hard $E^{\mathrm{par}}$ from $v_{\text{win}}$ onto the block-level
        $(u, a)$, honoured through $\iota$. The exemplar's own windows satisfy it
        trivially; proposals that break it are masked.
\end{itemize}

\paragraph{What the block level is, then.} Three channels on $\Lambda_B$:
\begin{center}
\small
\begin{tabular}{@{}llll@{}}
\toprule
channel & domain & reads & role \\
\midrule
$u$ & $\Z_m^2$ & $E_c$, $E_{\mathrm{par}}$, seam $\kappa$, $\nu$, $\mathrm{cost}(u, v)$ & the plan: appearance, objects, openings \\
$a$ & $\{\varnothing\} \cup K \times \Phi$ & overlap, contact, $\nu$, plan mask & the realisation: whole objects, hard \\
$v$ & small & compat, refine, $\mathrm{cost}(u, v)$ & the contract: reachability to the root \\
\bottomrule
\end{tabular}
\end{center}
with the update rates ordered $u \ll v \le a$: the plan is re-matched rarely
and expensively, the realisation is swept often and cheaply. The one thing
single $a$-moves could not do, merging, is no longer asked of them.

\paragraph{Where $\nu$ sits now.} $a$ may deviate from $\alpha^{\mathrm{own}}(u)$ at
cost $\nu$ (a different room of the same family, a shifted phase), but the plan's
openings are hard on $a$. So $\nu$ is a style term and $\iota$ a structural one.
Whether $\iota$ should also be $\nu$-soft (letting $a$ close a planned opening when
nothing fits) is the honourability question for this level: is there always a
realisation of a plan drawn from exemplar windows with matched interfaces? For
windows copied verbatim yes; for $a \ne \alpha(u)$ it must be checked.

%-----------------------------------------------------------------------------
\section{The interface, second pass: channels, views, factors}
\label{sec:views}

Sections~\ref{sec:objects}--\ref{sec:plan} make the coordinate look like a
promise: a value chosen above that the level below must honour. That is right,
and it is the simplification. The difference between a promise and a coordinate
is not in the channel but in how its statements are used. This section replaces
the member list of Section~\ref{sec:interface} with a smaller one.

\paragraph{Views.} A channel $c$ exposes a finite set of \emph{views}
\[
  \alpha_{c,k} : \Dom_c \to V_{c,k}, \qquad k \in \mathrm{views}(c),
\]
each a small-valued statement about the cell. The identity is a view. Examples:
\begin{center}
\small
\begin{tabular}{@{}lll@{}}
\toprule
channel & views & where they come from \\
\midrule
tile $t$ & port byte per side $\iota_d(t)$; node, exit, solid bits & tile set (ports \emph{are} views) \\
promise $v$ & identity; interface $\iota_d(v)$ & the language \\
object $a$ & identity; footprint bitmap; slots per side & templates \\
coordinate $u$ & $N_E(u)$; $\alpha^{\mathrm{own}}(u)$; $\iota_d(u)$; $\alpha_v(u)$ & \textbf{tables over exemplar positions} \\
certificate $(g,d)$ & identity & -- \\
\bottomrule
\end{tabular}
\end{center}
A coordinate differs from every other channel in one way only: its views are
read from precomputed tables indexed by exemplar position, so it exposes as
many statements as one cares to tabulate, and every tuple of them is
\emph{jointly realisable by construction}, since the exemplar is a real map. A
coordinate is a library of consistent promise bundles indexed by position.

\paragraph{Factors.} Every factor is a table over views, with offsets:
\[
  E\big(\alpha_{c_1,k_1}(z_{c_1}(p + o_1)),\; \dots,\; \alpha_{c_n,k_n}(z_{c_n}(p + o_n))\big),
  \qquad E : \textstyle\prod_i V_{c_i,k_i} \to \R \cup \{+\infty\},
\]
with $o_i \in \Z^2$ a same-level offset (the neighbour kinds), $o_i = 0$ (cell
kind) or $o_i$ pointing to the parent lattice (parent kind). Hard is a property
of the table ($\{0, +\infty\}$), not of the channel. Arity is usually 2; the
plan constraint $\big(a_p, a_q, \iota_d(u_p)\big)$ and the certificate's validity
$\big(t_p, g_p, d_p, t_q, g_q, d_q\big)$ are 3- and 6-ary, read as ``the other
arguments are fixed during this update''. A table may be stored
($V \times V$ floats, or bit planes) or computed (footprint bitmaps ANDed under a
shift); the update does not care which. Coupling through views keeps every
table small: $(\alpha^{\mathrm{own}}(u), a)$ is $|A| \times |A|$, never $m^2 \times |A|$.

\paragraph{Update.} A cell update on a block $B$ of channels enumerates (or
proposes) joint values, computes their views, sums every factor touching $B$
with the other arguments fixed, and draws. Exact when $\prod_{c \in B}|\Dom_c|$
is enumerable; bit-sliced when every factor on $B$ is a per-side relation on
port-group views; proposal-based otherwise, with the candidate set from
neighbours' values and the parent. Rates per block are free.

\paragraph{What ``promise'' and ``coordinate'' mean now.} Both are channels.
\begin{itemize}
  \item A \textbf{promise} is a channel some of whose views are honoured
        \emph{hard} by a lower level ($E^{\mathrm{par}}$ or $E^{\mathrm{cell}}$ tables
        in $\{0, \infty\}$), whose seam factor on its interfaces is exact (a
        compatible grid is realisable), and which has a hard parent relation
        $\refine$ (compositional $\alpha$).
  \item A \textbf{coordinate} is a channel whose views are honoured by a
        \emph{mix}: $\iota_d(u)$ hard on $a$ (the plan), $\alpha^{\mathrm{own}}(u)$ soft on
        $a$ ($\nu$), features soft on neighbours ($E_c$) and on tiles ($\tau$),
        $\alpha_v(u)$ soft on $v$ ($\mathrm{cost}$); whose seams are soft ($\kappa$)
        except on sides a promise requires; and whose parent relation is soft
        ($E_{\mathrm{par}}$, finite), so it has no $\refine$ and can never be stuck
        between levels.
\end{itemize}
So $u$ is a promise with respect to exactly its hard views, and realisability
of a hard view is witnessed per window by the exemplar and across seams by the
seam rule. The three properties a promise language needs (soundness,
realisability, exact seam rule) are asked of the \emph{hard views and their
tables}, view by view, not of the channel.

\paragraph{The data a level needs.} Enough to implement without further
decisions:
\begin{itemize}
  \item channels: $(h_c, \Dom_c, \{\alpha_{c,k}\})$, each view as an array over
        $\Dom_c$ (for $u$: arrays over exemplar positions, one per view);
  \item factors: $(\text{arguments } (c_i, k_i, o_i)_i,\ \text{table or kernel},\ \text{hard?})$;
  \item update blocks: a partition of $\Chan_h$ into jointly updated groups,
        each with its update kind, candidate generator if proposal-based, colouring
        (from the largest offset among its factors) and rate;
  \item per-level schedule: sweeps per block, ramp of any soft-to-hard weights,
        and the top level's prior.
\end{itemize}
The compile-time checks are then mechanical: symmetry of neighbour tables,
colouring spacing from offsets, never-empty for hard tables under the chosen
colouring, and a cost report $\sum_B |\Cand_B| \cdot \#\{\text{factors on } B\}$.

\paragraph{Mapping the existing code.} \texttt{port\_gibbs} is this interface at
$h = 1$ with views = port groups and factors = per-side relations. The promise
levels (\texttt{promise.py}) are identity-view channels with hard seam and parent
tables. \texttt{JointLevel} is a two-channel block $(u, v)$ with one cell table.
\texttt{macroobj} is a one-channel level ($a$) with a computed neighbour kernel
and a parent table $\nu$ against a fixed $u$. None of them is wrong; they are
four hand-specialised instances of one update.

%-----------------------------------------------------------------------------
\section{Promises are hard factors, not a kind of variable}
\label{sec:hard}

Section~\ref{sec:views} still names two usage patterns. They collapse further:
there are only channels and factors, and the hard factors are what used to be
called promises.

\paragraph{Hard and soft subsystems.} Split the factor set of a level into
$\mathcal F^{\mathrm{hard}}_h$ (tables in $\{0, +\infty\}$) and
$\mathcal F^{\mathrm{soft}}_h$ (finite). Write
\[
  H_h(z_{2h}) = \big\{ z_h : E(z_h, z_{2h}) = 0 \ \text{for every } E \in \mathcal F^{\mathrm{hard}}_h \big\}
\]
for the hard set of a level given its parent. The level conditional of
Eq.~\eqref{eq:level} is
\[
  g_h(z_h \mid z_{2h}) = \frac{\1[z_h \in H_h(z_{2h})]\; e^{-E^{\mathrm{soft}}_h(z_h, z_{2h})/T}}{Z_h(z_{2h})} .
\]
The hard subsystem fixes the support; the soft subsystem fixes the
distribution on it. The two have different sources and different obligations:
\begin{center}
\small
\begin{tabular}{@{}lp{2.9in}p{2.1in}@{}}
\toprule
 & hard factors & soft factors \\
\midrule
source & authored (sockets, templates, masks) or derived (compiled from a level below) & learned (counts, pseudo-likelihood, self-play) \\
count & a handful per channel type: seam, refine, honour & as many as declared \\
specific to & a channel \emph{type} (support, ports, objects), taking the tile set as data & an exemplar \\
checked by & the compiler (Section~\ref{sec:views}) & nothing; errors show as atypical output \\
tempered & never & by $T$ and the ramps \\
\bottomrule
\end{tabular}
\end{center}
This is the constitution's line: general rules, no rules specific to one tile
set. A hard rule belongs to a channel type and takes the tile set as data.

\paragraph{The three hard patterns.} Every hard factor so far is one of:
\begin{enumerate}
  \item \textbf{seam}: $(c, \iota_d, 0)$ with $(c, \iota_{\bar d}, e_d)$ at the same level
        (compat; the tile rules; overlap and contact on $a$);
  \item \textbf{refine}: $(c, \mathrm{id}, 0)$ with $(c, \mathrm{id}, \mathrm{par})$, the
        channel against its own coarser copy (merge consistency);
  \item \textbf{honour}: $(c, \alpha_k, \mathrm{par})$ with $(c', \mathrm{id}, 0)$ for a finer
        channel $c'$, or at one level $(c, \alpha_k, 0)$ with $(c', \mathrm{id}, 0)$
        (a view of a coarser or sibling channel equals what the finer one
        realises: $t = \mathrm{tmpl}[o]$, the plan mask on $a$, $A_b(x) = \pi_b$,
        validity of a certificate).
\end{enumerate}
A channel ``is a promise'' exactly when it appears in a honour factor as the
coarser or viewed argument. By that test $u$ (through $\iota$), $a$ (through
$\mathrm{tmpl}$), $v$, the object code $o$ and the certificate fields all are;
nothing is a promise by type.

\paragraph{Strong downward, soft elsewhere.} The user's description, read in
this notation: a channel's \emph{downward} factors (honour) are hard, its
\emph{sideways} factors (seam) are hard where the guarantee needs them and soft
otherwise, and its \emph{upward} factor (refine) is hard for derived channels
and soft ($E_{\mathrm{par}}$) for free ones. A guarantee $G$ holds for the output
iff it is implied by the hard subsystem alone, along an unbroken chain of
honour and seam factors from its certificate down to the tiles. A soft link
anywhere in that chain turns the guarantee into a tendency; it does not break
the sampler. That is the whole rule for deciding what is hard: make hard
exactly the factors on the chain of each guarantee, and nothing else.

\paragraph{What the correctness conditions attach to.} They are conditions on
$\mathcal F^{\mathrm{hard}}$ as a system, not on channels:
\begin{itemize}
  \item \textbf{soundness}: $z \in \bigcap_h H_h \Rightarrow G(x)$, by the chain argument above;
  \item \textbf{honourability}: $H_h(z_{2h}) \ne \varnothing$ for every $z_{2h}$ the levels
        above can produce, decided within a local halo (top-down never stuck);
  \item \textbf{never empty}: under the chosen colouring and update blocks the
        current value stays in every hard mask (local Gibbs never stuck);
  \item \textbf{exactness of seams} (quality, not correctness): a seam table no
        stronger than the tile-level rules imply, or coverage is lost.
\end{itemize}
The first three are what the compiler checks; the fourth is what the
self-play or the derived-rule compiler would eventually improve. For one
language these are the four properties of \texttt{coarse\_to\_fine.tex}
\S5.1; here they are asked once of the union of all hard tables at a level,
which is where joint honourability (support with connectivity, objects with
the plan) lives.

\paragraph{Coupling strength as a knob.} With hard and soft as two ends, a
factor may be ramped ($\lambda$ in the field sampler, $\mathrm{lam} \to \mathrm{HARD}$
in the macro runs) and ends hard; a run that ends with a finite weight has
turned a guarantee into a tendency for that run. Ramps are a sampler device
for mixing and do not change which factors are hard in the model.

%-----------------------------------------------------------------------------
\section{Promoting a view to a channel}
\label{sec:promote}

A view $\alpha_{c,k}$ is a deterministic function of $z_c$, so a hard factor on
it is a mask on $z_c$'s candidates. That is fine when $z_c$ is small and
enumerated (tiles: the port bytes are deterministic views and the tile update
scores every tile). It is a problem when $z_c$ is large and proposal-updated
(a coordinate): the hard factor rejects most candidates, and the exemplar may
hold no window that matches both the texture and the hard requirement, so the
channel freezes or coverage is lost. The remedy is the $(a, u)$ split of
Section~\ref{sec:objects}, as a general operation.

\paragraph{The operation.} Given a channel $c$ with view $\alpha : \Dom_c \to V$,
add a channel $w$ with $\Dom_w = V$ at the same level and
\begin{itemize}
  \item move every factor that read $\alpha(z_c)$ to read $w$ instead, hard ones
        included (seam, refine, honour by the level below);
  \item add one soft cell factor $E^{\mathrm{cell}}(w, z_c) = \mathrm{cost}(w, \alpha(z_c))$,
        either $\nu\,\1[w \ne \alpha(z_c)]$ or the calibrated form
        $-\log p(w \mid u)$, a smoothed histogram of $\alpha$ over exemplar windows
        near $u$ (\texttt{joint\_sampler.md}'s $\mathrm{cost}(u, v)$);
  \item give $w$ an exact update (small domain) and let $z_c$'s proposal update
        score $\mathrm{cost}(w, \alpha(\cdot))$ per candidate, or marginalise $w$ inside
        each candidate's score, $\bar\pi(u) \propto e^{-\Delta E(u)}\sum_w e^{-\mathrm{cost}(w, \alpha(u))}\1[w \text{ allowed}]$,
        which also credits a candidate for how many hard-feasible $w$ it admits.
\end{itemize}
As $\nu \to \infty$ the pair collapses to the deterministic view; the marginal
over $w$ of the promoted model is the original model with the hard factors
softened by exactly the cost table. So promotion is not a different model but
the same one with one knob, and the knob is what lets a hard system be
satisfied when the library has no exact match.

\paragraph{What a promise is, finally.} A promise in the old sense is a promoted
view of the coordinate: $v$ with $\mathrm{cost}(u, v)$ is $\alpha_v$ promoted, $a$
with $\nu$ is $\alpha^{\mathrm{own}}$ promoted. The pattern exists twice in the code
already. The rule for new hard factors follows:
\begin{quote}
  Put every hard factor on a small-domain channel. If the natural statement is a
  view of a large channel, promote the view. Keep a view deterministic only
  when its channel is enumerated exactly.
\end{quote}
Promotion costs one channel, one cost table and the masking of the large
channel's statement where the promoted one overrides it (Section~\ref{sec:objects});
it buys fast exact sweeps for the hard system, never-empty by construction,
and free update rates between the hard and the soft side.

\paragraph{Levels.} A promoted channel has a coarser copy where a refine factor
is wanted ($w_{2h}$ with a merge), and the coarser copy is itself the promoted
view of $u_{2h}$ with its own cost table. Honour runs downward: the level below
realises $w_h$ hard, not $\alpha(u_h)$. So the hard chain of a guarantee runs
entirely through promoted channels, and the coordinates touch it only through
cost tables. This is the picture the question asked for: hard variables, softly
coupled to the block channel, hard-coupled sideways and downward.

%-----------------------------------------------------------------------------
\section{Kernels: each update block owns its sampler}
\label{sec:kernels}

Hard channels are swept far more often than soft ones (they search; a hard
move is mostly rejected, so attempts must be cheap), and the fast samplers are
specialised: bit-sliced site draws, chain-block FFBS, closed-form interval
draws, kNN proposals. The generic layer should not try to be any of them. It
should own the factors and the schedule, and let each update block bring its
own kernel.

\paragraph{The contract.} An update block $B \subseteq \Chan_h$ comes with a
Markov kernel $K_B$ on $z_B$ (one cell, a colour class, or a chain of cells)
that leaves the conditional
\[
  \pi(z_B \mid z_{-B}) \;\propto\; \exp\!\Big(-\tfrac1T \sum_{E \ni B} E(\ldots)\Big)
\]
invariant, where the sum runs over \emph{every} factor with an argument in $B$,
the other arguments fixed. That is the whole requirement. Then any schedule
that applies kernels in sequence,
$K = K_{B_1}^{n_1} K_{B_2}^{n_2} \cdots$, leaves the level's joint invariant,
whatever the rates $n_B$, and is ergodic if the union of the blocks covers the
level and each kernel is irreducible on its conditional. A kernel that runs
many internal sweeps is the same as a high rate.

\paragraph{What a kernel may and may not do.}
\begin{itemize}
  \item It owns the \emph{storage} of its channels (tiles as port bytes and bit
        planes; rooms as label, phase and an occupancy grid; coordinates as
        integer pairs). Other channels never read that storage, only the views
        the channel publishes (Section~\ref{sec:views}).
  \item It must consume every factor touching $B$, including the cross-channel
        ones (cost tables from $u$, honour masks from a promoted channel, the
        plan mask). A factor it cannot consume exactly it may consume
        approximately and correct by Metropolis--Hastings (\texttt{port\_gibbs.md}
        \S5.3), which is still invariant.
  \item It may use any internal device whose correctness is its own business:
        a rotating partition of a large domain (the partition rule), a
        state-dependent candidate set (with the MH reverse set), a closed form
        over an interval, a chain block drawn by FFBS.
  \item It does not change other channels, and it does not ramp. Ramps and
        annealing belong to the schedule, as explicit weights on soft factors.
\end{itemize}

\paragraph{What the generic layer hands it.} The one real design decision: the
\emph{format of a factor restricted to $B$}. With the other arguments fixed, a
factor is a function $e : \Dom_B \to \R \cup \{+\infty\}$; kernels consume it in
different shapes, and a factor must be deliverable in the shape its consumer
wants:
\begin{center}
\small
\begin{tabular}{@{}lll@{}}
\toprule
kernel & consumes $e$ as & existing code \\
\midrule
enumerated site draw & a dense vector over $\Dom_B$ & \texttt{hier.pick}, macroobj A, \texttt{JointLevel} \\
bit-sliced site draw & quantised level planes plus a forbid mask & \texttt{bg\_sweep\_pairs}, \texttt{bg\_sweep\_field} \\
proposal (MH) & a callable per candidate & \texttt{\_mh\_S}, macroobj B \\
chain block (FFBS) & site vectors $e_i$ and transition tables $t_i$ & \texttt{envelope\_var.ffbs} \\
closed-form field & interval bounds from neighbours' values & \texttt{blockfield} \\
\bottomrule
\end{tabular}
\end{center}
A stored table converts to any of these (a level plane is a quantised table
row; a callable is a lookup). A computed factor (footprint overlap, the
dependant test) delivers a mask or a callable. So each factor declares the
shapes it can produce, each kernel the shape it consumes, and the compiler
checks that every factor on $B$ has a shape $K_B$ accepts. The cross-channel
terms a bit-sliced kernel sees from $u$ are then ``constant levels on masks'',
which \texttt{bg\_sweep\_field} already takes for its parent term; the
quantisation error is either accepted (a truncated target) or removed by the
MH correction.

\paragraph{Rates as the search budget.} With $u$ frozen, the hard channels'
kernels sample the conditional of the hard system given the soft one; the cost
tables are the only thing the search sees of the soft side. Sweeping the hard
system $n$ times per $u$-update is the budget for that search, and the free-rate
rule says any $n$ is valid. What $n$ affects is mixing: too small and the hard
system never equilibrates to the plan; too large wastes work and, with a
plan that keeps moving, chases a slowly moving target
(\texttt{macro\_objects.md} \S1). It is a knob of the schedule, reported next to
the per-kernel cost.

\paragraph{Division of labour, summarised.}
\begin{center}
\small
\begin{tabular}{@{}ll@{}}
\toprule
the channel brings & the generic layer owns \\
\midrule
storage, views, a kernel with its consumed shape & factors, with offsets and producible shapes \\
(domain, candidate generator if proposal-based) & update blocks, colouring, rates, ramps \\
 & restriction of factors to a block, in the kernel's shape \\
 & compile-time checks and the cost report \\
\bottomrule
\end{tabular}
\end{center}

%-----------------------------------------------------------------------------
\section{Composition, and what the compiler decides}
\label{sec:compiler}

\paragraph{Exact-kernel levels.} A Perlin-style terrain generator is a channel
too: its domain is a height (and slope, biome, \ldots) per cell, its kernel is
an exact draw from position-hashed noise, and it has no within-level factors.
In the chain graph it is a top component sampled once, never revisited, and
the levels below see it only through parent-kind cost tables on its views
(height $\to$ the support promise's cost; biome $\to$ a gate on the exemplar
library, the style channel of \texttt{dsl.md} \S7). The same holds for a split
grammar emitting regions, or a Wilson tree over windows: a kernel that samples
its conditional exactly in one pass needs no sweeps. What makes a cheap
generator and an MCMC level composable is the directed structure between
levels: local normalisation means the castle machinery never pushes back on the
terrain, so the terrain's kernel is exact regardless of what sits below. The
price is the one already paid (no bottom-up feedback), and the gain is that
levels are combined by views and cost tables only, with neither side knowing
the other's kernel.

\paragraph{What bounds the complexity.} Four things, all already in the
interface:
\begin{enumerate}
  \item the \textbf{DAG}: between levels directed, so the global schedule is a
        sequence of per-level schedules, each a cyclic order over a few blocks;
  \item \textbf{views}: a channel is known to others only by its published
        statements, so adding a generator adds view tables, not kernel knowledge;
  \item \textbf{declared couplings}: no table exists unless written, and the cost
        report prices each one;
  \item \textbf{bounded joint blocks}: channels tied by a hard cell factor must be
        updated jointly; the product domain is reported and capped.
\end{enumerate}
The explosion, when it comes, comes from soft couplings ($\binom{n}{2}$ tables)
and joint blocks (product domains). Both are visible at compile time.

\paragraph{Mechanical: structure and correctness.} From the declarations alone
the compiler can decide
\begin{itemize}
  \item the \textbf{order of levels}: a topological sort of the chain components by
        parent-kind factors; a partially directed cycle is an error (promises
        pointing both ways, \texttt{coarse\_to\_fine.tex} \S2);
  \item the \textbf{update blocks}: the connected components of channels under hard
        cell factors, which must be joint (else never-empty fails), with every
        other channel its own block; a cap on $\prod |\Dom|$ per block, with a
        warning and a suggested promotion when it is exceeded;
  \item \textbf{shape matching}: every factor on a block deliverable in its kernel's
        shape (Section~\ref{sec:kernels});
  \item the \textbf{colouring} of each block from the largest offset among its
        factors, plus one for dependant reads;
  \item the \textbf{hard-system checks}: symmetry of seam tables, never-empty under
        the colouring, honourability by enumeration on small halos, and
        optionally seam exactness against the level below;
  \item the \textbf{cost report}: per block per sweep,
        $|\Cand_B| \cdot \#\{\text{factors on } B\}$ in the kernel's units, and
        per level.
\end{itemize}

\paragraph{Not mechanical: rates, sweeps, ramps.} How many sweeps a block
needs, how often a soft channel is re-drawn per hard sweep, and how a hard
factor is ramped in are mixing questions, and mixing times are not computable
from the specification. Three ways to set them, in order of preference:
\begin{enumerate}
  \item \textbf{calibration}: run the level on finite windows with a generous
        schedule, measure per channel the autocorrelation of its views across
        sweeps and the decay of hard-factor violations under the ramp, and fit
        the smallest $n_B$ and ramp length that reach the plateau; emit a fixed
        schedule. This is an offline step like fitting the tables, and the
        output is a constant, so the fixed-step and locality rules hold at run
        time;
  \item \textbf{heuristics} as the starting point for calibration: exact-kernel
        levels one pass; small hard blocks many cheap sweeps; large soft blocks
        few expensive ones; rates in the ratio of their per-sweep costs so each
        block gets a comparable share of time;
  \item \textbf{declared} by the author, when the calibration is not worth running.
\end{enumerate}
Not allowed: deciding at run time from global statistics (a convergence test
over the map). A run-time rule must be local and bounded, as the ramps are,
and is then part of the schedule, not of the model.

\paragraph{The schedule as an object.} Per level: an ordered list of
$(B, n_B, \text{ramp}_B)$ repeated $s_h$ times, then the next level. It is small,
printable, and the only thing that changes between a calibrated run and a
default one. Everything correctness-related is fixed before it is chosen; the
schedule can make a sampler slow or poorly mixed but never invalid.

%-----------------------------------------------------------------------------
\section{Direction, and the free energy of the level below}
\label{sec:direction}

Take the simplest case: a block channel $\rho$ (dirt density) and the tiles
under it, with the honour factor $\mathrm{count}(x|_p) = f(\rho_p)$. In the
\emph{target} this factor is undirected: it constrains tiles given $\rho$ and
$\rho$ given tiles. In the \emph{generator} it is used top-down only. The
question is whether that loses anything.

\paragraph{The two models differ by a free energy.} Write $E_{h/2}(z_{h/2} \mid z_h)$
for the energy of the level below given the level above (its soft factors,
with the hard ones as masks). Its partition function given $z_h$,
\[
  Z_{h/2}(z_h) = \sum_{z_{h/2}} e^{-E_{h/2}(z_{h/2} \mid z_h)/T},
  \qquad
  F_{h/2}(z_h) = -T \log Z_{h/2}(z_h),
\]
is the \emph{free energy of the level below} as a function of the level above.
In one undirected energy over all levels, the marginal of $z_h$ is tilted by
$e^{-F_{h/2}(z_h)/T}$: a coarse value is weighted by how many fine
completions it admits and how good they are. In the chain graph it is not.
So ``both ways'' and ``top-down'' are different distributions on $z_h$, and the
difference is exactly $F_{h/2}$. This is not about mixing. A both-ways model
sampled perfectly and a top-down model sampled perfectly give different
coarse marginals.

\paragraph{Which one is right depends on where the coarse tables come from.}
\begin{itemize}
  \item If $z_h$ is \textbf{derived} and its tables are fitted as the
        \emph{push-forward} of the target, $\hat q(z_h) \approx q(z_h) = \sum_{x:\,\alpha(x) = z_h} p_G(x)$,
        then $-\log q = F + \text{const}$: the tilt is already in the tables,
        because a corpus of fine samples shows each coarse value as often as its
        completions warrant. Top-down is then exact (the $\mathrm{TV}$ bound of
        \texttt{coarse\_to\_fine.tex} \S3), and adding a bottom-up tilt would
        count $F$ twice.
  \item If $z_h$ is \textbf{designed} (Perlin heights, a grammar, a density the
        author draws), its distribution is a specification, and the tilt would
        corrupt it: heights with more tile completions would become more
        likely. Top-down is what the author means. A designed level sits above
        the learned ones and replaces their coarse marginal by a choice; the
        learned levels below realise it (the constitution's ``coarse levels are
        learned'' applies to the levels that model the exemplar, not to a
        specification placed above them).
  \item If $z_h$ is \textbf{free} ($u$), it is an auxiliary variable; its tables
        are the exemplar's statistics, and the pull-back toward $p$ happens in
        the tile-level relaxation, as now.
\end{itemize}
So, A: top-down only, as a rule of the model and not of the sampler. The
bottom-up direction is not a factor; it is the requirement that each level's
within-level energy be an approximation to $F$ of the level below, or be
intended as a specification that overrides it.

\paragraph{Where top-down breaks: the support of $F$.} Two objects that fully
determine the tiles under them cannot overlap, because $Z_{\text{tiles}}(a) = 0$,
$F = \infty$. A joint Gibbs on $(a, x)$ sees this with no extra rule: a move of
$a_p$ onto $a_q$'s footprint fails the honour factor against the \emph{current}
tiles, locally and immediately. A top-down pass cannot: when $a$ is sampled,
$x$ does not exist yet. The alternation the question mentions is the price of
the undirected model, and it breaks the single pass and chunked generation. The
top-down model must therefore carry $F$'s infinite part at the coarse level
itself, as the derived rule of Section~\ref{sec:plan}. This gives the general
statement:
\begin{quote}
  A top-down level is sound only if its hard factors contain $\mathrm{supp}\,F_{h/2}$,
  the set of coarse values with a feasible completion. Honourability is the
  same condition read downward.
\end{quote}
The checker that enumerates small halos for honourability is computing
$\mathrm{supp}\,F$; the derived overlap rule is $\mathrm{supp}\,F$ written by hand.

\paragraph{B: the learnable interaction terms are the finite part of $F$.}
$F_{h/2}(z_h)$ is a global function, but it is approximately local: the fine
level's energy is a sum of local factors, so $Z$ factorises over regions up to
boundary terms, and
\[
  F_{h/2}(z_h) \;\approx\; \sum_p f_1(z_h(p)) \;+\; \sum_{p \sim q} f_2\big(z_h(p), z_h(q)\big) \;+\; \cdots,
\]
a cluster expansion in which $f_1$ is the free energy of one cell's
completions, $f_2(z, z') = F_{\{p,q\}} - F_{\{p\}} - F_{\{q\}}$ the excess from two
neighbours having to agree across their seam, and higher terms are dropped. The
\emph{learnable interaction term} of the question is $f_2$ (and $f_1$), the
soft within-level tables of a coarse level, with three estimators:
\begin{center}
\small
\begin{tabular}{@{}lp{2.5in}p{1.8in}@{}}
\toprule
estimator & what it computes & where it fits \\
\midrule
derived (closed form, checker) & $\mathrm{supp}\,F$: the $\{0, \infty\}$ part & hard tables \\
counted from a corpus & $-\log \hat q$, i.e.\ $F$ up to a constant, by frequency & soft tables (counts, pseudo-likelihood) \\
self-play & $F_W(z_h|_W)$ on small windows $W$ by clamped fine-level runs & soft tables where no corpus exists \\
\bottomrule
\end{tabular}
\end{center}
Self-play, concretely: for each value pair $(z, z')$ of a two-cell window (and a
few halo values), clamp the coarse level, run the fine kernel, and estimate
$\log Z_W$ (thermodynamic integration or annealed importance sampling; or,
cheaply, the mean energy reached, which is the $T \to 0$ limit of $F$). Infeasible
runs give $f_2 = \infty$ and so recover the derived rule without writing it.
Cost: $|V|^2 \times$ halo values fine-level runs per table, which is why coupling
through views ($|V|$ small) and not through $u$ ($m^2$) is required.

\paragraph{The density example, closed.} With $\rho$ derived and counted from
the exemplar, $f_1(\rho)$ already contains $-\log \#\{\text{tile configurations
with that count}\}$ weighted by their energies: the entropy of each density
is in the table, and top-down is exact. With $\rho$ drawn from noise, the
author has chosen the density distribution and the tiles realise it; the
count honour factor is hard, nothing flows up, and that is the intent. In
neither case is a bottom-up factor written.

%-----------------------------------------------------------------------------
\section{The ideal, and the chain as its approximation}
\label{sec:ideal}

Interactions between levels will conflict in ways nobody wrote down. The
position taken here: the \emph{ideal} is the undirected joint of every declared
factor over every level, sampled by an algorithm too expensive to run
(Gibbs up and down all levels to equilibrium); the directed chain is its
approximation; and the learned terms are \emph{defined} as what makes the
approximation exact, with self-play as their estimator and cheaper things used
in practice. This section makes the three parts precise.

\paragraph{The ideal.} With levels $h_{\text{top}}, \dots, 1$ and the energies
$E_h(z_h \mid z_{2h})$ of Section~\ref{sec:direction} (within-level and parent
factors of level $h$, hard ones as masks),
\[
  p^*(z) \;=\; \frac{1}{Z^*}\, \exp\!\Big(-\tfrac1T \sum_h E_h(z_h \mid z_{2h})\Big).
\]
Every unforeseen conflict is in $p^*$: a coarse configuration whose fine
completions are all bad or all infeasible has small or zero marginal under
$p^*$, whether or not anyone anticipated it.

\paragraph{The identity.} Define, from the bottom up, the cumulative free
energy of everything below a level,
\[
  \tilde F_{1}(z_2) = -T\log \sum_{x} e^{-E_1(x \mid z_2)/T},
  \qquad
  \tilde F_{h}(z_{2h}) = -T\log \sum_{z_h} e^{-[E_h(z_h \mid z_{2h}) + \tilde F_{h/2}(z_h)]/T}.
\]
Then, with the corrected level energies $\tilde E_h = E_h + \tilde F_{h/2}$,
\begin{equation}
  p^*(z) \;=\; \prod_h \tilde g_h(z_h \mid z_{2h}),
  \qquad
  \tilde g_h(z_h \mid z_{2h}) = \frac{e^{-\tilde E_h(z_h \mid z_{2h})/T}}{\sum_{z_h'} e^{-\tilde E_h(z_h' \mid z_{2h})/T}} ,
  \label{eq:telescope}
\end{equation}
by telescoping: each normaliser is the partition function of everything below
$2h$ given $z_{2h}$, which is $e^{-\tilde F_h(z_{2h})/T}$, and the next level up
cancels it. So \emph{a directed chain whose level energies include the
cumulative free energy of the levels below is the undirected joint, sampled
top-down in one pass}. The chain as used, $\tilde F \equiv 0$, is the
approximation; the only thing it drops is $\tilde F$. With estimates $\hat F$ in
place of $\tilde F$,
\[
  \log p^*(z) - \log g(z) \;=\; \sum_h \big(\hat F_{h/2} - \tilde F_{h/2}\big)(z_h) \;+\; \text{(the same errors, re-normalised one level up)},
\]
so every discrepancy between the ideal and the chain is an error in some
$\hat F$, and an error at level $h$ shows up as a wrong marginal at level $2h$,
nowhere else.

\paragraph{The learned terms, defined.} The soft within-level tables of a
coarse level are an approximation to $\tilde F_{h/2}$ in cluster-expanded
form (Section~\ref{sec:direction}): a per-cell term, a pair term per
adjacency, both over views. This is the definition the question asks for: a
learned coupling between channels at level $h$ is a term of the cluster
expansion of the free energy of the corrected level $h/2$. It has three
properties worth having.
\begin{itemize}
  \item It is \textbf{recursive and two-level}. $\tilde F_h$ is computed from
        $E_h$ and $\tilde F_{h/2}$ only, so self-play at level $h$ runs the
        \emph{corrected} kernel of level $h/2$, not the whole stack below. Fit
        from the bottom: tiles first, then each level against the one below it.
  \item It \textbf{absorbs the unforeseen}. Any conflict among lower factors,
        hard or soft, local or carried by a certificate, raises $\tilde F$ on the
        coarse values that provoke it. Nothing has to be anticipated; it has
        to be measured.
  \item It is \textbf{the sampler's own free energy if measured with the sampler}.
        Running the fixed-step fine kernel under clamping, rather than an exact
        sampler, estimates the free energy \emph{as realised by that kernel at
        that budget}: the coarse level then prefers values the real fine stage
        can complete well, which is what cross-scale consistency as a training
        goal means (constitution, non-goals).
\end{itemize}

\paragraph{Self-play as the estimator.} For a view pair $(v, v')$ on two
adjacent cells and a small set of halo values: clamp the coarse channel, run
the corrected kernel of the level below for its scheduled sweeps, record
\[
  \hat F_W \approx -T\log Z_W \ \ \text{(thermodynamic integration or AIS)}
  \qquad\text{or}\qquad
  \hat F_W = \langle E \rangle \ \ (\text{the } T \to 0 \text{ proxy}),
\]
and set $f_2(v, v') = \hat F_{\{p,q\}} - \hat F_{\{p\}} - \hat F_{\{q\}}$, $f_1(v) = \hat F_{\{p\}}$
up to a constant, $+\infty$ on runs that end infeasible. Cost:
$|V|^2 \times$ halos $\times$ one fine run per pair table, feasible because
views are small.

\paragraph{Cheaper in practice, in order of fidelity.}
\begin{enumerate}
  \item \textbf{A corpus.} If fine samples exist (an exemplar, an exact reference
        sampler), counting the push-forward gives $-\log \hat q = \hat F + \text{const}$
        with no runs at all; this is self-play someone already did.
        Pseudo-likelihood fits the same thing as conditionals.
  \item \textbf{Mean energy.} $\langle E \rangle$ instead of $-T \log Z$: drops the
        entropy of the completions; right at $T \to 0$, biased toward coarse
        values with few good completions over many fair ones.
  \item \textbf{Support only.} The honourability checker: $\hat F \in \{0, \infty\}$.
        This is the derived hard rules, and it is the minimum a sound level
        needs (Section~\ref{sec:direction}).
  \item \textbf{Truncation.} Pairs along adjacency only, halos marginalised by
        a few draws, views coarser than the channel.
  \item \textbf{Nothing}: channels couple only through tiles and hard rules,
        $\hat F \equiv 0$ beyond the support. The current state of the code.
\end{enumerate}
Each step down is a known omission from the identity, so the price of each
is nameable: the entropy term, the finite part, the off-adjacency terms.

\paragraph{In practice.} Exact sampling of $p^*$ is not the goal; plausible,
coherent output under parameters the author controls is. The identity is
still useful under that goal, for two reasons. It says which omissions are
safe: the support must be right (a wrong support is a stuck sampler or an
invalid map), the finite part only shifts statistics, and a shifted statistic
is a knob. And it says what a knob \emph{is}: a weight on a soft table is a
tempered estimate of a free energy, so turning it up or down has a meaning (more
or less credit for how well the level below can complete a value) rather than
being a tuning accident. The working default is therefore: support from the
checker or by hand; pair and unary tables from counts where a corpus exists and
from the mean-energy proxy where it does not; each table with a weight the
author sets. Refit by self-play only where a run shows coherence failing across
a level boundary.

\paragraph{What self-play cannot see, and what fixes it.} The cluster
expansion is local; a conflict that is nonlocal at the lower level
(reachability across many cells) does not appear in a pair term. It appears
once a certificate or promise has made it local, as a term on the certificate's
views. So certificates and self-play are complementary: certificates make
$\tilde F$ local, self-play measures the local $\tilde F$. The rule for an
unforeseen conflict discovered in a run is then: if it is local, it will be
absorbed by refitting; if it is not, it needs a certificate first.

%-----------------------------------------------------------------------------
\section{Worked example: a tree with mass-monotone branches}
\label{sec:tree}

Guarantee wanted: every branch block is joined to a branch block of equal or
higher mass. What has to be written, in the interface of
Sections~\ref{sec:views}--\ref{sec:hard}.

\paragraph{The channel that carries branches.} Whatever block-level channel $z$
holds the tree: segments as objects $a$ (mass from the template), a coordinate
$u$ (mass tabulated over exemplar positions), or tiles aggregated per block
(mass as $\alpha$ of the block's wood cells). The recipe does not depend on
the choice; only where the view tables come from.

\paragraph{Views (two, plus one bit).} Deterministic if $z$ is enumerated
exactly (objects, aggregated tiles); promoted to channels with a cost table if
$z$ is a coordinate (Section~\ref{sec:promote}). Everything below reads the same
either way.
\begin{itemize}
  \item $\mathrm{mass}(z) \in \{0, 1, \dots, M\}$, with $0$ meaning no branch;
  \item $\mathrm{join}_d(z) \in \{0,1\}^{B}$, the cells on side $d$ at which wood
        reaches the seam (an interface, as $\iota_d$ for rooms). Two adjacent
        blocks are \emph{joined} when
        $\mathrm{join}_d(z_p) \wedge \mathrm{join}_{\bar d}(z_q) \ne 0$;
  \item $\mathrm{root}(z) \in \{0,1\}$: a trunk block, which needs no witness.
\end{itemize}

\paragraph{The rule as stated is local but not sound.} ``Some joined neighbour has
mass $\ge$ mine'' is a hard factor of radius 1 and can be used as is. But two
blocks of equal mass joined to each other and to nothing else satisfy it, so
equal-mass islands float, and with strict $>$ the walk ends at a local maximum
that need not be a trunk. The guarantee that is probably meant, \emph{every branch
block is joined by a chain of non-decreasing mass to a trunk}, is the
certificate pattern: add a rank.

\paragraph{The certificate channel.} $d_p \in \{0, \dots, D\} \cup \{\infty\}$ at the
block level, small, updated jointly with $z_p$. Validity:
\[
  V_p \;=\; [\mathrm{mass}(z_p) = 0] \;\lor\; \mathrm{root}(z_p) \;\lor\;
  \exists\, q \in N_1(p):\ \mathrm{joined}(p,q)\ \land\ \mathrm{mass}(z_q) \ge \mathrm{mass}(z_p)\ \land\ d_q < d_p ,
\]
with $\mathrm{root}(z_p) \Rightarrow d_p = 0$ and $\mathrm{mass} = 0 \Rightarrow d_p = \infty$. One hard
factor, $L \cdot \1[\lnot V_p]$, reading $(z_p, d_p)$ and $(z_q, d_q)_{q \in N_1(p)}$:
arity 10, radius 1. This is \texttt{blockfield}'s $(g, d)$ field with the gap
count $g$ dropped and the mass comparison added.

\paragraph{Soundness (three lines).} Follow witnesses from any branch block:
$d$ strictly decreases, so the walk ends, and it can end only at a root; every
step is a join, and mass never decreases along it. So every branch block is
joined by a non-decreasing chain to a trunk. The check is local: count invalid
blocks.

\paragraph{The update.} Block $p$ redraws $(z_p, d_p)$ given its neighbours:
\begin{enumerate}
  \item for each candidate $z$: compute $\mathrm{mass}, \mathrm{join}_d, \mathrm{root}$ (table
        reads), and for each neighbour $q$ whether $q$ is a possible witness
        ($\mathrm{joined}$, $\mathrm{mass}_q \ge \mathrm{mass}$) and whether $q$ is a
        \emph{dependant} (its only current witness is $p$; computed once per update
        from the neighbours' own neighbourhoods);
  \item the valid $d$ for a fixed $z$ form an interval: $d > \min_{\text{witnesses}} d_q$
        from own validity, $d < \min_{\text{dependants}} d_q$ from theirs, $d = 0$ if
        root, $d = \infty$ if mass $0$ and no dependant; a dependant that would lose
        its join or its mass bound under $z$ forbids $z$ outright;
  \item weight: the soft factors on $z$ (exemplar terms, $\nu$, $\kappa$) plus an
        optional $\delta\, d$ pulling toward short routes, summed over the interval
        in closed form; draw $z$, then $d$ within its interval.
\end{enumerate}
Never empty: the current $(z_p, d_p)$ satisfies every mask, as long as the
dependant test is reproduced exactly. Colouring: $V_q$ for a neighbour reads
$q$'s neighbours, so same-colour blocks must be $\ge 3$ apart: the 5-colouring.

\paragraph{Infinite worlds.} A chain to a trunk crosses chunk boundaries, which
the block level may not follow. So a coarser level must commit, as
\texttt{macro\_full} does for rooms: each window either contains a trunk or
names a \emph{root side} toward its parent window; a block on that side may take
a virtual witness across it (mass $M$, $d = -1$). That is one refine factor
(the window tree) and one honour factor (the virtual witness), both already
written for rooms. Without it the rule still holds within a chunk, but a tree
cannot be guaranteed to reach a trunk that lies elsewhere.

\paragraph{Mixing.} Hard from the start, a branch can only grow outward from a
trunk one block per sweep and only shrink from the tips, a ratchet; the field
sampler's experience is that a soft ramp ($\lambda\, \1[\lnot V_p]$ rising to $L$)
lets the soft factors lay out branches first and the certificate then prunes
and reconnects. The ramp does not change the model's hard set.

\paragraph{What to write, as code.}
\begin{center}
\small
\begin{tabular}{@{}ll@{}}
\toprule
piece & size \\
\midrule
three view tables over $\Dom_z$ (mass, join per side, root) & precompute \\
the certificate channel $d$ and its interval draw & reuse of \texttt{blockfield}, minus $g$ \\
the validity kernel $V$ and the dependant test & one function, $\sim$50 lines \\
the window-level root side (if infinite) & reuse of \texttt{macro\_full}'s tree \\
compiler entries: arity, radius, colouring 5, never-empty & declarations \\
\bottomrule
\end{tabular}
\end{center}
No table is learned; nothing is specific to the tree's exemplar.

%-----------------------------------------------------------------------------
\section{Adding a channel: the recipe and its cost}
\label{sec:recipe}

Suppose many channels exist and one more is to be added, with its own
constraints. What has to be written, and what the machinery does.

\paragraph{The steps.}
\begin{enumerate}
  \item \textbf{Level and domain.} Pick $h_c$ and $\Dom_c$, and whether the channel is
        derived (a function of the tiles under a cell), designed (a generator)
        or free.
  \item \textbf{Views.} Write each $\alpha_{c,k}$ once, as a function of a cell's
        tiles (derived) or of the value (designed, free). The machinery applies
        it to the exemplar to make the view tables, and for a coordinate at the
        same level the cost table $-\log p(\alpha_{c,k} \mid u)$ by counting.
  \item \textbf{Hard rules.} Write the seam predicate on the interface view (if the
        constraint is sideways), the honour predicate (how the level below
        realises the value: a mask on the lower channel's views), and the merge
        (if the channel has a coarser copy). Each is a small table or a
        predicate of a few lines.
  \item \textbf{Kernel.} Small domain: the enumerated site draw, nothing to write.
        Large domain: promote the views (Section~\ref{sec:promote}) and write a
        candidate generator for the large part. A nonlocal constraint: add a
        certificate channel and reuse the field kernel with a new validity
        predicate.
  \item \textbf{Couplings.} Declare any soft coupling to another channel beyond the
        defaults. The default is none; the automatic ones are the channel's own
        unary, pair and parent tables and its cost table to the coordinate, all
        counted.
  \item \textbf{Compile.} The checks, colouring, block membership, default schedule
        and cost report run without input. Calibrate rates if the default is
        visibly wrong.
\end{enumerate}

\paragraph{Written versus automatic.}
\begin{center}
\small
\begin{tabular}{@{}lp{2.2in}p{2.2in}@{}}
\toprule
item & written per channel & automatic \\
\midrule
view functions & 1--4 small functions & tables over $\Dom$ or exemplar positions \\
unary, pair, parent tables & -- & counted from the exemplar through the views \\
cost table to $u$ & -- & counted (the promoted-view coupling) \\
seam, honour, merge & 1--3 predicates or tables & conversion to each kernel's shape \\
couplings to other channels & 0 by default; one table each if declared & counted or self-played if asked \\
kernel & 0 (enumerated); 1 generator (large); 1 predicate (certificate) & restriction of factors, draw, MH bookkeeping \\
colouring, blocks, schedule & -- & from offsets, hard cell factors, defaults \\
checks & -- & symmetry, never-empty, honourability, cost \\
tests & the predicates against enumeration on small cases & invariants of the generic layer \\
\bottomrule
\end{tabular}
\end{center}

\paragraph{Three sizes.}
\begin{itemize}
  \item \textbf{A soft derived channel} (a material class per block that should
        follow the exemplar): one view function; no hard rule; everything else
        counted. Tens of lines. It changes nothing elsewhere.
  \item \textbf{A hard local channel} (a density that tiles must match exactly, with
        a seam rule): one view, one honour mask on tiles, one seam table. The
        tile kernel receives one more factor in its shape and is otherwise
        untouched. Under a hundred lines.
  \item \textbf{A hard nonlocal channel} (the tree of Section~\ref{sec:tree}): the
        views, a certificate channel, one validity predicate with its dependant
        test, and a window-level commitment reused from the rooms. A few
        hundred lines, most of it the predicate and its test.
\end{itemize}
For comparison, each channel today is a module with its own sweep, energy,
state and experiment script: \texttt{macroobj} 1050 lines, \texttt{macrocontact}
560, \texttt{blockfield} 560, \texttt{generic} 1050.

\paragraph{What is paid once.} The recipe is cheap only because the generic
layer exists: view tables and counting, factor restriction in each kernel's
shape, the enumerated and proposal kernels, the field kernel with pluggable
validity, masking of the coordinate's statement under a promoted channel, the
checker, the schedule compiler and the cost report. That is the refactor, and
it is where the lines go. Of it, the pieces not in the code in any form are
factor restriction with shapes, the checker, the schedule compiler and the
masking; the rest is consolidation.

\paragraph{The condition behind the recipe.} The counts above assume kernels
take a stack of factors and never a named term, so that a channel acting on
tiles adds an entry to the tile kernel's stack and not a line to its code.
Section~\ref{sec:stacking} states that rule, the closure it needs, and the one
addition to the factor interface (gates) that keeps two channels' statements
about the same tiles from being counted twice.

\paragraph{What nobody escapes.} Two costs stay per channel whatever the
machinery. A hard rule must be \emph{right}: sound, with an exact seam, and
the checker only finds the second kind of error. And an existing channel that
should \emph{respond} to the new one needs a declared coupling, which is a
decision, not code; the default of zero is what keeps $n$ channels from
costing $n^2$ tables, and it is also why a new channel is invisible to the old
ones until someone says otherwise.

%-----------------------------------------------------------------------------
\section{Many channels on one kernel}
\label{sec:stacking}

Two channels both change the energy of the tiles under them. If the tile
kernel was written to take \emph{a} parent term, the second channel forces a
rewrite, and the third again. The recipe of Section~\ref{sec:recipe} holds only
if this cannot happen. What prevents it is a rule on the kernel side of the
contract, a closure property of the shapes, and one addition to the factor
interface.

\paragraph{The rule: a kernel takes a list, never a named term.} The input to a
kernel for block $B$ at site $p$ is the \emph{stack} of restricted factors,
$\{e_1, \dots, e_K\}$, each in the kernel's shape, from whatever channels
declared them. The kernel sums (or ANDs) the stack and draws. It never knows
which channel a factor came from, how many there are, or at what level they
live. Adding a channel that affects tiles adds one entry to the stack. This
is what \texttt{bitgibbs.c} already does with \texttt{acc\_add} over factor
stacks; the generic version makes it the only way a kernel receives energy.
A kernel written against a fixed neighbourhood (``the four sides and the
parent'') violates the rule; one written against ``the factors on $B$, each
with its offsets'' satisfies it.

\paragraph{Closure: every shape must add.} For the stack to be summed inside
the kernel, each shape needs a sum (soft) and a meet (hard):
\begin{center}
\small
\begin{tabular}{@{}lll@{}}
\toprule
shape & soft sum & hard meet \\
\midrule
dense vector over $\Dom_B$ & vector add & mask AND \\
level planes + forbid plane & bit-sliced adder, within a level budget & forbid OR \\
callable per candidate & sum of calls & conjunction \\
chain tables $e_i, t_i$ & add site and transition tables & mask AND on both \\
interval bounds & -- & intersect intervals \\
\bottomrule
\end{tabular}
\end{center}
One real cost hides here: the bit-sliced kernel has a \emph{level budget}
($\mathrm{AB} = 5$ bits, 32 levels). Each soft factor declares its level range,
and the compiler checks the sum fits; if it does not, it coarsens the unit
(losing resolution, corrected by MH) or moves a factor out of the planes into
the MH correction. So in that kernel a new soft channel on tiles has a price in
bits, and the compiler reports it next to the time cost.

\paragraph{Any coarser level may speak to the tiles.} A parent-kind factor's
offset may point to any ancestor lattice, not only $2h$: a block channel at
$h = 8$ and a window channel at $h = 64$ both restrict to a tile site by
reading their cell above it. The chain graph is unchanged (level $h$ is
conditioned on everything above it, and the order is still a topological
sort); the restriction just reads the ancestor at the right level.

\paragraph{Double counting, and masks as part of a factor.} Summing is the
right default for independent statements (an object's template and a density
count), and wrong for two statements \emph{about the same thing}: the coordinate's
texture term and an object's template both say what the tiles under a
footprint look like, and summing counts the evidence twice (\texttt{dsl.md}
\S6). This is a modelling fact, not a kernel problem, and it is expressed in
the factor, not in the kernel: a factor may carry a \textbf{gate}, a view of
some channel at some offset, and contributes only where the gate is open,
\[
  E(\ldots)\cdot \1\big[\gamma(z_{c'}(p + o'))\big].
\]
The coordinate's tile term is gated off under any footprint; the object's
honour mask is gated on there. Precedence among channels is then authored, as
a small ordered list, and compiled into gates. The kernel still sees a stack;
some entries are zero at some sites. For a view that \emph{is} the gate
(ownership), promotion already gives it a channel to live on.

\paragraph{Joint satisfiability of hard stacks.} Two hard honour masks on one
tile site AND together, and the meet can be empty: an object's template cell
and a density count that cannot both hold. This is joint honourability at the
tile level, and it is the checker's job (the union of a level's hard tables,
Section~\ref{sec:hard}), not the kernel's: the kernel only requires
never-empty, which holds when the state is consistent. A new channel whose
honour mask can conflict with an existing one is caught at compile time, on
small halos, before any sweep.

\paragraph{Who knows what.} Three layers, each party seeing one:
\begin{center}
\small
\begin{tabular}{@{}lll@{}}
\toprule
layer & written by & known to \\
\midrule
views of a channel & the channel's author & anyone who couples to it \\
tables over views (factors) & the author connecting two channels & the generic layer \\
shapes (a kernel's input format) & the kernel's author & the generic layer \\
\bottomrule
\end{tabular}
\end{center}
An upper channel that constrains tiles writes a table over (its view, a tile
view); it knows the tile channel's \emph{views} (port bytes, node bit, solid
bit) and nothing about the tile kernel's shape, other channels, or how many
of them act on tiles. The tile kernel knows its shape and nothing about
levels or channels. The generic layer converts tables to shapes and stacks
them. So the ``communication channel'' of the question is the pair
(views published downward, tables written on them), and the format a kernel
consumes is hidden behind the conversion. The only crossing of the line is a
\emph{computed} factor (footprint overlap, a dependant test), which cannot be a
stored table and must produce a shape directly: its author writes a producer
for the shapes it supports, and the compiler refuses a kernel it does not.
Those are few, and they are the hard structural rules of a channel type,
written once with it.

\paragraph{What still forces a kernel change.} Exactly two things. A factor in
a shape the kernel cannot consume (a radius-2 soft pair term into the
per-side bit-sliced kernel): then the compiler routes it to the MH correction
or to the enumerated kernel, and the author may choose to extend the kernel
instead. And a new structural device (a certificate whose validity does not
fit the pluggable predicate of the field kernel, a block update of a new
shape): then it is a new kernel, under the same contract, for that channel's
block only. Neither touches an existing kernel.

%-----------------------------------------------------------------------------
\section{Open questions}
\label{sec:open}

\paragraph{Settled in this document (by argument, not experiment).}
\begin{itemize}
  \item Objects: ownership form on a block-level coordinate, with the object as
        a promoted view $a$ coupled by $\nu$ (Sections~\ref{sec:objects},
        \ref{sec:promote}). Interlocking is a different overlap table.
  \item Connectivity: planned in the coordinate update through interfaces,
        honoured hard by $a$, global half as the window promise
        (Section~\ref{sec:plan}). The forest heuristic is replaced, not fixed.
  \item Interface: channels expose views; factors are tables on views with
        offsets; hard is a property of the table; a promise is a channel in a
        honour factor (Sections~\ref{sec:views}, \ref{sec:hard}).
  \item Hard factors live on small channels; promote a view when it does not
        (Section~\ref{sec:promote}).
  \item Each update block brings its own kernel under one invariance contract;
        the generic layer delivers restricted factors in the kernel's shape
        (Section~\ref{sec:kernels}).
  \item Top-down only, as a rule of the model; soundness requires the hard
        factors to contain $\mathrm{supp}\,F$ of the level below
        (Section~\ref{sec:direction}).
  \item Learned terms are the cluster expansion of the cumulative free energy
        below, with self-play as the definition and counts, mean energy or
        support as the practical estimators (Section~\ref{sec:ideal}).
\end{itemize}

\paragraph{Still open.}
\begin{enumerate}
  \item \textbf{The experiment.} Contact-mode macro at full scale with the
        $(u, a)$ split and the interface seam factor: does planning openings in
        $u$ remove the patchwork, and does $a \ne \alpha(u)$ ever leave a planned
        opening unrealisable (the honourability question of Section~\ref{sec:plan})?
  \item \textbf{The proposal set for $u$ given $a$.} Instances of $(k, \phi)$ in
        the exemplar as a precomputed list, plus coherent shifts, parent and $k$
        links; the MH reverse set where it depends on the current value.
  \item \textbf{Factor shapes.} The concrete representation of a restricted
        factor for each kernel (dense vector, level planes, callable, chain
        tables, interval bounds) and the conversions between them; whether
        quantised cross-channel terms into the bit-sliced kernel are corrected
        by MH or accepted as truncation.
  \item \textbf{Joint blocks.} The cap on a product domain and what the compiler
        suggests when it is exceeded (promotion, a chain block, a partition
        schedule).
  \item \textbf{Calibration.} The measurement that sets sweeps and rates per
        block (view autocorrelation, violation decay under a ramp) and how it is
        run on finite windows.
  \item \textbf{The honourability checker.} Enumeration on small halos over the
        union of a level's hard tables; its cost as hard channels multiply, and
        whether per-channel checks ever suffice.
  \item \textbf{Self-play tooling.} Clamped two-cell runs of a corrected kernel
        with the mean-energy proxy; when a refit is worth running.
  \item \textbf{A coarser commitment to objects.} Whether a window-level cell
        must promise objects it does not own (an envelope, a count), or whether
        $\varnothing$ always being admissible plus the window promise is enough.
\end{enumerate}

\end{document}
